\documentclass[preprint,12pt]{elsarticle}
\usepackage[T1]{fontenc}
\usepackage[utf8]{inputenc}
\usepackage[brazil]{babel}
\usepackage{amsmath,amssymb,booktabs,graphicx}
\usepackage{xurl}
\usepackage{placeins}
\setlength{\emergencystretch}{2em}
\renewcommand{\floatpagefraction}{0.8}
\renewcommand{\topfraction}{0.9}
\renewcommand{\textfraction}{0.1}
\usepackage[hidelinks]{hyperref}
\journal{Supply Chain Analytics -- versão de trabalho em português}
\begin{document}
\begin{frontmatter}
\title{Seleção de políticas de reposição por perfil de demanda: avaliação temporal em uma rede varejista}
\author[ufal]{Matheus Inácio Batista Santos}
\affiliation[ufal]{organization={Instituto de Computação, Universidade Federal de Alagoas},country={Brasil}}
\begin{abstract}
A escolha de políticas por perfil de demanda acrescenta complexidade à reposição, mas seu valor depende do desempenho em períodos posteriores. Comparamos essa estratégia com a seleção de uma política global em 145 séries loja--produto de uma rede varejista brasileira. Sete candidatas são avaliadas em duas origens temporais, com treino, validação e teste separados e três sementes nos métodos estocásticos, totalizando 4.350 ajustes. Cada seletor minimiza custo entre candidatas com atendimento médio de pelo menos 0,70 na validação das demais lojas; a avaliação usa o teste da loja reservada. A seleção por perfil apresenta custo médio 1,93\% maior e atendimento 3,89 pontos percentuais superior. Os intervalos condicionais de 95\% são, respectivamente, $[-0{,}32;4{,}16]\%$ e $[2{,}78;5{,}07]$ pontos percentuais. Os casos abaixo da meta diminuem de 119 para 100 entre 283 casos série--origem com demanda positiva; sete casos têm atendimento indefinido. A diferença concentra-se na primeira origem; na segunda, os seletores fazem escolhas idênticas. O resultado explicita um compromisso entre custo e serviço, sem estabelecer economia geral. A avaliação combina diferenciação de políticas e restrição por grupo, sem isolar o efeito da representação. Coorte retrospectiva, janelas curtas parcialmente sobrepostas e inferência condicional limitam a generalização.
\end{abstract}
\begin{keyword}
Inventory control \sep Policy selection \sep Demand heterogeneity \sep Service level \sep Retail
\end{keyword}
\end{frontmatter}

\section{Introdução}
Uma rede varejista precisa decidir quando e quanto repor para muitas combinações de loja e produto. Uma regra que funciona para vendas frequentes pode produzir excesso de estoque em séries esparsas; outra, econômica em média, pode não cobrir picos relevantes. O problema não se resume a prever demanda: envolve escolher uma regra de decisão e ajustar seus parâmetros sob custos e exigências de atendimento \cite{Axsater2006InventoryControl}.

Comparações de algoritmos ampliam o conjunto de alternativas disponíveis. Políticas analíticas, heurísticas, busca por parâmetros e aprendizado por reforço oferecem diferentes compromissos entre custo, serviço e esforço computacional. Entretanto, a existência de um portfólio não resolve a escolha operacional. O gestor precisa saber se vale aplicar uma política uniforme ou distinguir grupos de séries, e se a distinção oferece ganho suficiente para justificar maior complexidade.

Formulamos a escolha como um problema de seleção de algoritmos: características de uma instância ajudam a escolher um método de um conjunto candidato \cite{rice1976,Kotthoff2014AlgorithmSelection}. Em inventário, a instância é uma série loja--produto, e a recomendação depende do custo e do serviço exigido. A referência global aplica a mesma família de política a todas as séries, com parâmetros ajustados individualmente. O seletor por perfil pode mudar essa família entre grupos. Portanto, a comparação mede o valor da diferenciação da escolha, mantendo a calibração por série em ambas as estratégias.

Estudos recentes relacionam previsão e desempenho operacional \cite{sca_f81ad2b3e023}, comparam algoritmos de controle \cite{sca_27e994c54848} e examinam a informação decisória preservada por critérios de classificação \cite{sca_c5c06dcbc9e5}. Esses enfoques motivam avaliar o que a seleção por perfil acrescenta ao portfólio: distinguir padrões de demanda só justifica uma regra mais complexa se a distinção melhorar a decisão em relação a uma referência global competitiva.

O artigo concentra-se em uma pergunta: selecionar políticas por perfil de demanda melhora o compromisso entre custo e atendimento frente à seleção de uma política global? As duas estratégias usam a mesma informação passada e o mesmo portfólio. A evidência reportada corresponde à matriz completa da Bahia. O conjunto público M5 integra a avaliação complementar em execução e não fornece resultados quantitativos neste manuscrito.

Três análises sustentam essa pergunta: o compromisso custo--serviço das candidatas; a diferença entre seleção global e por perfil no teste; e a estabilidade das escolhas e do atendimento entre origens temporais. A comparação de controladores serve para interpretar a seleção, sem constituir uma segunda proposta de algoritmo. A contribuição é uma avaliação empírica do valor incremental dessa escolha em teste temporal, com perfis calculados no treino, exclusão da loja da seleção e análise da insuficiência de atendimento individual.

\section{Trabalhos relacionados e posicionamento}
A classificação da demanda organiza diferenças de frequência e tamanho das ocorrências. A taxonomia ADI--CV$^2$ é uma referência para distinguir demanda suave, errática, intermitente e \textit{lumpy} \cite{syntetos2005}. Uma classe estatística, porém, não identifica sozinha a melhor política: custo de falta, prazo, estado inicial e requisito de atendimento também influenciam a decisão.

Integrar classificação e controle já é um objetivo estabelecido. Mohammaditabar et al.\ \cite{Mohammaditabar2012Integrated} propõem escolher grupos e políticas conjuntamente, com solução por recozimento simulado. Sheikh-Zadeh et al.\ \cite{SheikhZadeh2021PBIC} desenvolvem classificação orientada ao desempenho para um sistema de múltiplos itens e escalões. Esses trabalhos impedem tratar a simples associação entre perfil e política como novidade. Nosso recorte é diferente: regras de perfil fixas e seleção de candidatas já ajustadas são confrontadas com uma referência global, usando validação em outras lojas e teste temporal na loja reservada. Não implementamos os métodos desses autores como comparadores e não alegamos superioridade sobre eles.

Tadayonrad e Ndiaye \cite{sca_f81ad2b3e023} aproximam a avaliação de previsão dos custos e do serviço de inventário. Essa ligação sustenta o uso de consequências operacionais como critério: neste estudo, precisão preditiva não substitui custo e atendimento. A política de estoque limitado CappedBaseStock acrescenta uma referência motivada pela literatura de vendas perdidas \cite{Xin2021CappedBaseStock}; propriedades teóricas dessa família não são garantias para nossa calibração e cenário.

Zabraoui et al.\ \cite{sca_27e994c54848} comparam famílias de algoritmos de inventário e apresentam combinações de algoritmos genéticos (GA) e aprendizado por reforço com dados M5. A arquitetura GA--DQN, que combina GA e uma rede Q profunda (DQN, do inglês deep Q-network), e a ideia de portfólio já estão, portanto, na literatura da revista. Nosso foco não é propor a mesma combinação como novidade, mas avaliar a camada de escolha entre alternativas. As heurísticas inspiradas naquele estudo envolvem decisões de implementação próprias, e custos, filtros e granularidade não constituem uma reprodução exata de seus experimentos.

A análise de critérios de Theunissen et al.\ \cite{sca_c5c06dcbc9e5} destaca a diferença entre importância estatística e utilidade decisória. Não implementamos seu método topológico. Testamos uma consequência próxima: mesmo um perfil interpretável pode ser pouco útil se recomendar praticamente a mesma política global ou se suas mudanças elevarem custo sem benefício de serviço relevante.

Revisões de aprendizado por reforço em inventário discutem a importância do ambiente, da recompensa e da avaliação \cite{Boute2022DRLInventoryRoadmap}. Uma baixa taxa de atendimento ou frequência de pedidos extrema não caracteriza, isoladamente, uma família algorítmica. Os resultados abaixo descrevem implementações e matrizes específicas. A interpretação exige identificar o controlador, seu orçamento e as condições comuns de avaliação.

\input{protocolo_bahia_20261015}

\section{Seleção e protocolo de comparação}
\subsection{Referência global e escolha por perfil}
Para um conjunto $A$ de casos série--origem usados na escolha, seja $A^+$ o subconjunto com demanda positiva no bloco correspondente. Definimos
\begin{equation}
 \bar C_{\pi,A}=\frac{1}{|A|}\sum_{i\in A}C_{i\pi},\qquad
 \bar F_{\pi,A}=\frac{1}{|A^+|}\sum_{i\in A^+}F_{i\pi}.
\end{equation}
Todos os casos entram no custo; somente os casos com atendimento definido entram em sua média. Quando $A^+$ é vazio, a média de atendimento não é definida. Na validação da Bahia não há denominadores nulos; no teste, eles são contabilizados separadamente. Cada caso definido recebe o mesmo peso, sem ponderação pelo volume demandado. Seja $\Pi$ o portfólio de políticas. A política global é
\begin{equation}
 \widehat\pi_A=\arg\min_{\pi\in\Pi:\bar F_{\pi,A}\geq\alpha}\bar C_{\pi,A}.
\end{equation}
A seleção por perfil aplica a mesma regra separadamente a $A_g$, o subconjunto associado ao perfil $g$. Se nenhuma política atende à restrição, escolhe-se maior atendimento médio, com custo como desempate, e registra-se fallback. Esse caso é explicitado: não é uma escolha viável que cumpriu o limiar. Entre candidatas viáveis com mesmo custo, preferimos maior serviço; no fallback empatado em serviço, menor custo. Empates restantes usam ordem alfabética do identificador. Um perfil sem informação válida de atendimento nas lojas de escolha utiliza a regra global e recebe marcação de contingência.

A restrição é média e específica do conjunto usado para escolher. Não exige $F_{i\pi}\geq\alpha$ em toda série. Também não garante a média da loja excluída. Além disso, exigir viabilidade em cada perfil é diferente de exigi-la apenas no agregado. Uma política global viável pode ser inviável em um grupo; a escolha por perfil não precisa apresentar menor custo mesmo na avaliação retrospectiva. A comparação avalia duas estratégias operacionais com diferentes formas de distribuir a exigência de serviço. Ela não isola o efeito dos descritores de demanda de um efeito de impor a restrição em cada grupo.

O limiar principal é $\alpha=0{,}70$, fixado no cenário experimental. Ele não representa recomendação de meta comercial. Avaliamos também 0,80, 0,85, 0,90 e 0,95 na etapa de seleção, sem retreinar o portfólio para cada meta. A calibração das candidatas mantém referência de 0,70; a sensibilidade descreve escolhas entre essas candidatas, não a melhor política possível para cada limiar. O aumento relativo de custo é
\begin{equation}
 \Delta_C=100\frac{\bar C_{\mathrm{perfil}}-\bar C_{\mathrm{global}}}{\bar C_{\mathrm{global}}}.
\end{equation}
$\Delta_C>0$ indica aumento; $\Delta_C<0$, redução. Diferenças de atendimento são dadas em pontos percentuais. Custo menor acompanhado por atendimento menor não é chamado de dominância.

\subsection{Perfis operacionais e informação disponível}
Os perfis são regras heurísticas ordenadas, não clusters aprendidos nem políticas previamente validadas. A prioridade é Sparse High Impact, High Volatility Seasonal, Unstable Trend, Low Volume Stable e Fast Moving. A primeira regra utiliza ADI elevado, CV$^2$ elevado e sequência longa de zeros; a segunda, autocorrelação sazonal ou concentração de ocorrências; a terceira, tendência ou entropia. Os perfis de baixo volume estável e venda frequente não aparecem nos prefixos avaliados da Bahia. Para tornar a regra reproduzível, usamos $a=n/\max(n_+,1)$, $c^2=[s/(\mu+10^{-9})]^2$ e $Z$, maior sequência de zeros. Aqui, $n$ é o número de ciclos do prefixo de treino, $n_+$ é a quantidade de ciclos com demanda positiva, $\mu$ é a média sobre todos os ciclos e $s$ é o desvio-padrão amostral, incluindo zeros. Essa variância difere da calculada apenas sobre os tamanhos positivos da demanda na taxonomia original. A Tabela~\ref{tab:profile_rules} apresenta os cortes, aplicados na ordem indicada; uma série sem correspondência recebe Sparse High Impact.

\begin{table}[htbp]\centering\small
\caption{Regras operacionais de perfil: primeiro corte satisfeito determina o grupo.}\label{tab:profile_rules}
\begin{tabular}{@{}p{3.7cm}p{7.2cm}@{}}\toprule
Perfil & Condição\\\midrule
Sparse High Impact & $a\geq1{,}32$, $c^2\geq0{,}49$ e $Z\geq5$\\
High Volatility Seasonal & Autocorrelação na defasagem 17 $\geq0{,}25$ ou burstiness $\geq0{,}30$\\
Unstable Trend & Módulo da tendência normalizada $\geq0{,}05$ ou entropia $\geq2$\\
Low Volume Stable & $\mu\leq2$ e $c^2<0{,}49$\\
Fast Moving & $a<1{,}32$\\\bottomrule\end{tabular}\end{table}

Burstiness é $(s_\Delta-\mu_\Delta)/(s_\Delta+\mu_\Delta+10^{-9})$, em que $\mu_\Delta$ e $s_\Delta$ são a média e o desvio-padrão dos intervalos $\Delta$ entre ocorrências positivas; recebe zero com menos de dois intervalos. Tendência é a inclinação estimada por mínimos quadrados ordinários (OLS) sobre o índice de ciclo, dividida por $\mu+10^{-9}$. Entropia é $-\sum_j p_j\log_2(p_j+10^{-12})$, em que $p_j$ é a frequência relativa na classe $j$ de um histograma com dez intervalos de mesma largura. Autocorrelação é a covariância na defasagem 17 dividida pela variância global; recebe zero se a série tem até 18 ciclos ou variância desprezível. Os rótulos descrevem regras locais, não sazonalidade ou impacto econômico empiricamente validados.

Na avaliação temporal, essas características são calculadas exclusivamente no prefixo de treino. A eventual associação de um nome de perfil a uma política no código não determina a escolha: custo e serviço na validação definem a recomendação.

Com 23 ciclos, a autocorrelação na defasagem 17 usa somente seis pares de observações; com 17 ciclos, ela não é estimada. A mudança de perfil entre origens pode, portanto, decorrer da disponibilidade e instabilidade dos descritores. Os limiares não foram ajustados pelo teste nem constituem uma classificação operacional validada externamente.

\subsection{Avaliação e incerteza}
As regras global e por perfil usam a validação das demais lojas e são aplicadas ao teste da loja reservada. Cada política é ajustada apenas no treino da respectiva série; o seletor usa resultados de validação, sem consultar o teste para escolher. As sementes são promediadas dentro de cada caso série--origem antes das comparações. Custos incluem todos os casos; médias de atendimento excluem os denominadores nulos, contabilizados separadamente. Origens são medidas repetidas, com um ciclo de teste compartilhado.

A reserva da loja se aplica à escolha da família de política. O controlador conserva o treino no histórico da própria série e o estado inicial usa seu passado observado. O desenho avalia a transferência da regra de escolha entre lojas com histórico disponível; não representa implantação em uma loja sem observações anteriores.

A incerteza das diferenças de custo e atendimento é descrita por intervalos percentis de 95\%, obtidos em 5.000 reamostragens pareadas das 115 lojas. Cada loja sorteada leva todas as suas séries e origens. Controladores e seletores permanecem fixos nos sorteios; os intervalos são condicionais às escolhas e aos treinamentos realizados. Para atendimento, as somas e os denominadores positivos são recalculados em cada sorteio. A dependência entre lojas que compartilham produtos permanece uma limitação.

\input{bahia_20261015}

\section{Discussão e implicações operacionais}
A avaliação temporal responde à pergunta central com um compromisso: os perfis aumentam o atendimento, mas também o custo. A melhora concentra-se na primeira origem, enquanto a segunda reproduz integralmente a escolha global. O aumento de custo é uma estimativa pontual cujo intervalo condicional inclui zero; isso não comprova igualdade de custos. O intervalo do atendimento permanece positivo sob a reamostragem adotada, sem estabelecer que o efeito se repetirá em novas trajetórias.

Essa interpretação é diferente de comparar apenas as políticas individuais. GA--DQN melhora custo e atendimento frente ao DQN, mas não frente ao GA no conjunto dos testes. Mesmo que um controlador fosse superior em média, ainda seria necessário demonstrar que as características permitem identificar quando recomendá-lo. A distinção entre qualidade dos candidatos e utilidade do seletor aproxima o problema da seleção de algoritmos \cite{rice1976,Kotthoff2014AlgorithmSelection}.

A comparação conceitual com Zabraoui et al.\ \cite{sca_27e994c54848} situa o portfólio e a arquitetura híbrida na literatura. Recortes, implementações e orçamentos distintos impedem confirmar ou refutar seu ranking com os valores deste estudo. O foco em consequências operacionais também se relaciona à avaliação de custo e serviço discutida por Tadayonrad e Ndiaye \cite{sca_f81ad2b3e023}. Precisão de previsão ou identidade do algoritmo não substituem essa comparação.

A agregação temporal exige atenção. Na primeira origem, ambas as estratégias ficam abaixo de 0,70 no teste; reunir as origens produz médias acima desse valor. O atendimento ponderado permanece inferior à média por série, indicando que o serviço não se distribui uniformemente pelo volume de demanda. Uma decisão empresarial precisa definir se a meta se aplica à média das séries, às unidades demandadas ou a itens críticos individualmente. Para o gestor, o resultado oferece uma decisão concreta: avaliar se o serviço adicional compensa o custo e o uso de regras diferenciadas. Essa avaliação exige calibrar custos e criticidade comercial; o limiar de 0,70 do cenário não deve ser convertido em meta operacional por conveniência.

\section{Limitações e continuidade}
O ajuste e os perfis usam apenas o passado da origem, mas a coorte mantém identificadores escolhidos retrospectivamente. Os testes cobrem sete ciclos e compartilham um ciclo entre origens. Séries da mesma loja e produtos presentes em várias lojas podem ser dependentes. A reamostragem por loja preserva parte dessa estrutura, sem representar novos treinamentos ou novas realizações da demanda.

As vendas podem omitir demanda censurada por falta de estoque. Custos são unidades do cenário e não demonstram economia monetária realizada. O estado inicial comum é produzido por uma política de referência; o teste não mede a continuação do estoque particular de cada candidata. Estoque terminal sem crédito, horizonte curto e valores de custos podem alterar o compromisso observado. Essas condições precisam acompanhar a interpretação dos resultados. Além disso, a calibração de CappedBaseStock incorpora penalidades de estabilidade que não são usadas no GA, e as candidatas não têm esforço computacional total equivalente. A matriz permite comparar o portfólio especificado, sem atribuir os contrastes apenas à família algorítmica.

O M5 é um conjunto público de vendas de varejo associado a uma competição de previsão \cite{Makridakis2022M5Accuracy}. Sua execução de reposição deverá avaliar a mesma pergunta, com preparação e elegibilidade documentadas e estimativas separadas por base. A execução em andamento ainda não demonstra generalização. A continuidade inclui mais períodos, menor sobreposição dos testes, sensibilidade econômica e restrições de serviço por série. Um seletor supervisionado é uma extensão possível, mas sua validação não é necessária para interpretar a comparação entre as duas regras explícitas aqui avaliadas.

\section{Conclusão}
Em 145 séries da Bahia, a seleção por perfil acrescenta 1,93\% ao custo e 3,89 pontos percentuais ao atendimento médio, reduzindo os casos abaixo de 0,70 de 119 para 100 entre 283 casos definidos. O efeito varia entre períodos e desaparece como diferenciação de escolha na segunda origem. A evidência sustenta avaliar o valor do serviço adicional, sem afirmar economia geral ou superioridade do seletor.

A contribuição é separar desempenho dos controladores, escolha da política e avaliação temporal da recomendação. O resultado permanece condicional à coorte, às janelas e ao cenário. A avaliação complementar no M5 e em outras janelas é necessária para examinar a extensão desse resultado a outros contextos. Diferenciar grupos só terá utilidade operacional se o benefício se mantiver frente a uma referência global competitiva e justificar seu custo.

\section*{Disponibilidade de dados e código}
Os dados transacionais brasileiros são proprietários e não são redistribuídos. Código, protocolo, scripts de análise e fontes do manuscrito estão disponíveis em \url{https://github.com/matheus013/msc}. Dados, bancos de execução e modelos não integram a versão atual desse repositório. Uma versão arquivada e identificável do pacote de reprodução, bem como os derivados cuja divulgação seja autorizada, deverão acompanhar a submissão. Não se promete acesso aos registros restritos mediante solicitação.
\section*{Autoria, financiamento e interesses}
Autoria, contribuições CRediT, financiamento, agradecimentos e declaração de interesses permanecem sujeitos à confirmação dos autores responsáveis. Esta versão de trabalho em português requer consolidação dessas informações e adequação ao idioma e às normas editoriais antes da submissão.
\section*{Uso de inteligência artificial na preparação do manuscrito}
OpenAI Codex foi utilizado para auxiliar organização bibliográfica, leitura de código, preparação de scripts de análise e visualização e redação. As figuras apresentam dados calculados pelos scripts, sem geração de imagens por modelo generativo. A versão atual exige verificação e revisão do autor humano responsável, que deverá conferir fontes, cálculos e afirmações e assumir responsabilidade pelo texto final. Sistemas de IA não são autores.

\bibliographystyle{elsarticle-num}
\bibliography{references_ptbr}
\end{document}
